% 03 — OR: Parallel switching circuit (Battery -> parallel switches p,q -> Lamp)
\begin{tikzpicture}[line width=1pt,>={Stealth[length=2.4mm]},font=\itshape]
  \ctikzset{bipoles/length=0.9cm}
  % Battery on left and lamp on right
  \draw[navy,line width=1.2pt] (0,0) to[battery2,l=$+ \quad -$] (0,3);
  \draw[navy,line width=1.2pt] (0,3) -- (1,3);
  \draw[navy,line width=1.2pt] (1,3) -- (1,1);
  % สวิตช์ p: หน้าสัมผัสวงกลมและคันโยกแบบเดียวกับภาพวงจรอนุกรม
  \draw[navy,line width=1.2pt] (1,3) -- (2,3);
  \draw[navy,line width=1.2pt,fill=white] (2,3) circle (0.09);
  \draw[navy,line width=1.2pt,fill=white] (3.7,3) circle (0.09);
  \draw[navy,line width=1.2pt] (2.08,3.04) -- (3.58,3.34);
  \draw[navy,line width=1.2pt] (3.7,3) -- (4.6,3);
  \node[font=\itshape\large] at (2.85,3.48) {$p$};
  % สวิตช์ q ในแขนงขนาน ใช้สัญลักษณ์เดียวกัน
  \draw[navy,line width=1.2pt] (1,1) -- (2,1);
  \draw[navy,line width=1.2pt,fill=white] (2,1) circle (0.09);
  \draw[navy,line width=1.2pt,fill=white] (3.7,1) circle (0.09);
  \draw[navy,line width=1.2pt] (2.08,1.04) -- (3.58,1.34);
  \draw[navy,line width=1.2pt] (3.7,1) -- (4.6,1);
  \node[font=\itshape\large] at (2.85,1.48) {$q$};
  \draw[navy,line width=1.2pt] (4.6,1) -- (4.6,3);
  \draw[navy,line width=1.2pt] (4.6,3) -- (5.6,3) to[lamp] (5.6,0);
  \draw[navy,line width=1.2pt] (5.6,0) -- (0,0);
  \node[font=\itshape\small,black] at (3,-0.6) {วงจรสวิตช์ไฟฟ้าแบบขนาน (OR Operation)};
\end{tikzpicture}
